\section{Литературное расширение explainability-графа}

\subsection{Роль дополнительного литературного слоя}

В предыдущей главе explainability-слой был описан как knowledge graph, который интерпретирует shortlist-кандидатов,
сформированный моделью CatBoost. Однако rule layer в базовой версии графа в основном опирается на сигналы, извлеченные
из train-only статистики переходов между культурами, групповой принадлежности и регионального контекста. Такой подход
хорошо согласуется с данными проекта, но оставляет открытым вопрос о том, можно ли дополнить граф \textbf{внешними
агрономическими источниками}, описывающими желательные и нежелательные переходы между культурами.

Для ответа на этот вопрос в работе был добавлен отдельный этап \texttt{12\_literature\_augmented\_graph.ipynb}. Его
задача состоит не в замене предиктивного ядра и не в построении новой экспертной системы, а в расширении уже
существующего explainability-слоя дополнительными literature-backed сигналами. Тем самым knowledge graph получает еще
один интерпретационный уровень: помимо data-driven переходов в нем появляются support- и warning-сигналы, опирающиеся на
русскоязычные агрономические источники.

Принципиально важно, что данный этап не изменяет финальную модель диплома. Как и ранее, финальным предиктивным ядром
остается baseline CatBoost, а литературные правила используются только как \textbf{дополнительный семантический
контекст} для анализа shortlist-кандидатов. Поэтому литературное расширение следует рассматривать как отдельную
explainability-надстройку, а не как новый механизм выбора следующей культуры.

\subsection{Источники данных для этапа}

Литературное расширение графа построено как отдельный воспроизводимый notebook-этап, который использует уже готовые
артефакты explainability-анализа и добавляет к ним curated-таблицу литературных правил. Входными данными этапа выступают
два класса источников.

Во-первых, используются артефакты базового explainability-этапа \texttt{09\_knowledge\_graph\_explainability}. К ним
относятся таблицы \texttt{case\_summary.csv}, \texttt{case\_explanations.csv}, \texttt{kg\_nodes.csv},
\texttt{kg\_edges.csv}, \texttt{kg\_rules.csv}, \texttt{transition\_edges.csv} и метаданные запуска
\texttt{run\_meta.json}. Именно они задают исходный граф, фиксированный набор case studies и базовые значения
\texttt{graph\_support\_score}, с которыми затем сравнивается literature-aware версия графа.

Во-вторых, используется curated-вход \texttt{literature\_transition\_rules.csv}, сохраненный в директории
\texttt{artifacts/results/literature\_graph/}. Эта таблица формируется вручную на основе отобранных русскоязычных
источников и содержит краткие правила вида \texttt{crop\_context} $\rightarrow$ \texttt{candidate}, тип сигнала
(\texttt{support} или \texttt{warning}), силу сигнала, уровень сопоставления (\texttt{direct} или
\texttt{group\_approx}), а также поле \texttt{zonal\_note} для случаев, где применимость правила ограничена
определенными агроклиматическими зонами.

Результаты этапа сохраняются в отдельную директорию \texttt{artifacts/results/literature\_graph/}. Основными итоговыми
таблицами являются \texttt{literature\_augmented\_case\_explanations.csv},
\texttt{literature\_augmented\_case\_summary.csv} и \texttt{literature\_vs\_base\_case\_comparison.csv}. Кроме того,
сохраняются визуализации literature-aware графов для тех же четырех case studies, которые были использованы в базовом
stage 09.

\subsection{Использованные русскоязычные источники}

В curated-таблицу были включены только те правила, которые можно связать с конкретными русскоязычными публикациями или
учебно-методическими материалами. Для широких переходов между зерновыми, зернобобовыми, паром и кормовыми культурами
использовались материалы по севооборотам и предшественникам культур
\cite{stgau-rotation-basics,wheat-predecessor-study,kubsau-soy-predecessors}. В частности, они позволили добавить
support-сигналы для переходов к \texttt{wheat} после \texttt{forage\_hay} и \texttt{fallow}, а также support- и
warning-сигналы для переходов \texttt{corn} $\rightarrow$ \texttt{soybeans}, \texttt{corn} $\rightarrow$
\texttt{legumes} и повторных или чрезмерно насыщенных зерновых звеньев.

Отдельно были использованы русскоязычные источники по хлопчатнику
\cite{cotton-irrigated-technology,cotton-repeat-crop-study,cotton-pest-uzbekistan}. Они позволили заменить часть прежних
приближенных \texttt{group\_approx}-сигналов на прямые правила для класса \texttt{cotton}. В текущей версии графа это
реализовано как:

\begin{itemize}
\tightlist
\item
  support для переходов \texttt{legumes} $\rightarrow$ \texttt{cotton};
\item
  support для переходов \texttt{forage\_hay} $\rightarrow$ \texttt{cotton};
\item
  warning для перехода \texttt{cotton} $\rightarrow$ \texttt{cotton};
\item
  zonal warning для узкого звена \texttt{wheat} $\rightarrow$ \texttt{cotton}.
\end{itemize}

При этом по хлопчатнику была сохранена принципиальная оговорка о зональности. В отличие от зерновых и зернобобовых
переходов, часть cotton-правил относится прежде всего к хлопководческим и орошаемым регионам. Поэтому в таблице правил и
в notebook сохраняется поле \texttt{zonal\_note}, а сами сигналы трактуются как \textbf{direct but zonal}: они прямые
для культуры \texttt{cotton}, но не претендуют на универсальную общероссийскую применимость.

\subsection{Способ интеграции литературных правил в граф}

Литературные сигналы не смешиваются с train-derived переходами неявным образом. Для каждого сработавшего literature-rule
в графе добавляется отдельный support- или warning-узел интерпретации, связанный с кандидатом из shortlist. Тем самым в
literature-aware графе сохраняется различие между двумя типами evidence:

\begin{itemize}
\tightlist
\item
  signals, наблюдаемыми в данных проекта;
\item
  signals, пришедшими из внешних русскоязычных источников.
\end{itemize}

На уровне candidate-level explainability это приводит к пересчету \texttt{graph\_support\_score}. Новый score
по-прежнему не трактуется как вероятность и не заменяет модельный ранг; он лишь суммирует более широкий набор
support/warning-сигналов. Визуально и содержательно это означает, что на тех же четырех кейсах stage 09 можно увидеть,
как меняется относительная graph-based поддержка отдельных кандидатов после добавления literature-backed правил.

\subsection{Результаты на тех же case studies}

Литературное расширение было применено к тем же четырем кейсам, что и в базовом explainability-анализе:
\texttt{case\_high\_confidence}, \texttt{case\_medium\_confidence}, \texttt{case\_low\_confidence} и
\texttt{case\_difficult\_class}. Это важно, поскольку позволяет сравнивать не разные выборки наблюдений, а один и тот же
набор ситуаций до и после добавления литературного слоя.

На уровне итоговой case-level категории согласованности заметного изменения не произошло. В
\texttt{case\_high\_confidence} статус \texttt{strong} сохранился, а для трех остальных кейсов сохранился статус
\texttt{weak}. Иными словами, литературный слой не изменил распределение итоговых alignment-категорий на рассматриваемых
case studies. Это означает, что формального улучшения итоговой метрики согласованности в данном ограниченном наборе
кейсов зафиксировано не было.

Сводное сравнение по тем же кейсам приведено в таблице~\ref{tab:literature-effect}. Она показывает, что литература
влияет прежде всего на интерпретационную поддержку отдельных кандидатов, а не на итоговую категорию согласованности.

\begin{table}[H]
\centering
\scriptsize
\caption{Изменение graph support после добавления литературных правил}
\label{tab:literature-effect}
\setlength{\tabcolsep}{2pt}
\begin{tabular}{@{}>{\raggedright\arraybackslash}p{0.13\textwidth}
                >{\raggedright\arraybackslash}p{0.12\textwidth}
                r
                r
                r
                >{\raggedright\arraybackslash}p{0.18\textwidth}
                >{\raggedright\arraybackslash}p{0.14\textwidth}@{}}
\toprule
Case & Top-1 модели & \makecell[r]{Base\\$gs$} & \makecell[r]{Lit\\$gs$} & \makecell[r]{$\Delta gs$\\top-1}
    & \makecell[l]{Лучший literature-aware\\кандидат} & Alignment \\
\midrule
High confidence
  & \texttt{forage\_hay}
  & 0.661
  & 0.591
  & $-0.070$
  & \texttt{forage\_hay} (0.591)
  & strong $\rightarrow$ strong \\
Medium confidence
  & \texttt{corn}
  & 0.376
  & 0.236
  & $-0.140$
  & \texttt{soybeans} (0.598)
  & weak $\rightarrow$ weak \\
Low confidence
  & \texttt{legumes}
  & 0.011
  & 0.071
  & $+0.060$
  & \texttt{soybeans} (0.698)
  & weak $\rightarrow$ weak \\
Difficult class
  & \texttt{corn}
  & 0.220
  & 0.220
  & $\phantom{+}0.000$
  & \texttt{soybeans} (0.386)
  & weak $\rightarrow$ weak \\
\bottomrule
\end{tabular}

\vspace{0.4em}
\footnotesize $gs$ --- \texttt{graph\_support\_score}; $\Delta gs$ --- разность между literature-aware и базовым score
для top-1 кандидата модели.
Источник значений: \path{artifacts/results/literature_graph/literature_augmented_case_summary.csv}.
\normalsize
\end{table}


Тем не менее на уровне отдельных кандидатов наблюдаются содержательные сдвиги. В \texttt{case\_medium\_confidence}
литература заметно ослабила повтор \texttt{corn}: его \texttt{graph\_support\_score} снизился с 0.376 до 0.236, тогда
как для \texttt{soybeans} score вырос с 0.508 до 0.598. В \texttt{case\_low\_confidence} top-1 кандидат \texttt{legumes}
получил умеренное усиление с 0.011 до 0.071, но наиболее сильную literature-aware поддержку получила альтернатива
\texttt{soybeans} со score 0.698. В \texttt{case\_difficult\_class} литература не изменила статус top-1 у \texttt{corn},
но снизила поддержку повторного \texttt{cotton} с 0.385 до 0.305, что согласуется с добавленным warning-правилом для
повторного размещения хлопчатника. Наконец, в \texttt{case\_high\_confidence} статус согласованности не изменился, хотя
повтор кормового звена был слегка ослаблен: score для \texttt{forage\_hay} снизился с 0.661 до 0.591.

Если рассмотреть усредненную динамику по всем candidate-level строкам comparison-таблицы, то средняя разница
\texttt{literature\_graph\_support\_score - base\_graph\_support\_score} оказывается положительной и составляет примерно
0.011. Однако для model top-1 кандидатов средняя дельта отрицательна и составляет около $-0.038$, тогда как для не-top-1
кандидатов она положительна и составляет около 0.035. Дополнительно показательно, что из 12 candidate-level строк 4
получили положительную дельту, 3 отрицательную и 5 остались практически без изменений. Это хорошо согласуется с
наблюдаемым поведением literature-aware слоя: он не столько усиливает уже выбранный top-1, сколько \textbf{корректирует
shortlist}, штрафуя повторы и поддерживая более разнообразные или агрономически осмысленные альтернативы.

\subsection{Интерпретация результатов и ограничения}

Полученный результат важен именно своей интерпретацией. Литературные источники в данной работе не дали измеримого
прироста итоговой case-level alignment-метрики на выбранных четырех кейсах. Поэтому было бы некорректно представлять
этот этап как улучшение predictive quality или как новое доказательство превосходства системы.

Вместе с тем литературный слой сыграл полезную роль как \textbf{semantic correction layer}. Он сделал candidate-level
объяснения более содержательными, позволил явным образом зафиксировать warning-сигналы для повторных или нежелательных
переходов и усилил некоторые агрономически правдоподобные альтернативы в shortlist. Особенно полезным это оказалось в
спорных кейсах, где top-1 кандидат модели не обязательно является единственным правдоподобным вариантом и где
пользователю важно видеть не один ответ, а набор конкурирующих интерпретируемых альтернатив.

Ограничения этапа также следует обозначить явно. Во-первых, сам набор литературных правил является curated и
относительно компактным. Во-вторых, часть правил построена на уровне \texttt{group\_approx}, поскольку агрегированные
классы проекта не всегда имеют прямое соответствие одной конкретной культуре из источника. В-третьих, для
\texttt{cotton} прямые правила доступны в основном в zonal-формулировке, поэтому их переносимость за пределы
хлопководческих и орошаемых регионов должна трактоваться осторожно. Наконец, анализ выполнялся на тех же четырех case
studies, а не на полном тестовом наборе, поэтому он описывает прежде всего \textbf{качество интерпретации}, а не новую
глобальную метрику всего пайплайна.

\subsection{Выводы по главе}

Добавление отдельного literature-aware этапа показало, что explainability-граф можно расширить внешними русскоязычными
агрономическими источниками без изменения финального предиктивного ядра. Такой слой сохраняет общую архитектурную логику
работы: CatBoost формирует shortlist, confidence-слой оценивает надежность результата, а knowledge graph интерпретирует
кандидатов. Литературные правила естественно встраиваются именно в последний из этих уровней.

Практический эффект этого расширения состоит не в пересмотре финального model ranking, а в более содержательной
интерпретации shortlist. На рассмотренных кейсах литературные сигналы в основном усиливали agronomically plausible
alternatives и штрафовали нежелательные повторные звенья, но не изменили итоговую категорию согласованности между model
top-1 и graph evidence. Поэтому литературное расширение корректнее всего трактовать как \textbf{дополнительную
explainability-главу}, усиливающую содержательность рекомендаций, но не превращающую прототип в полноценную экспертную
систему или универсальный оптимизатор севооборота.
